% Lösungen Punkte und Strecken im Koordinatensystem · Körper und Drehungen · Prüfungs-Fokus Abitur GK (zusammenbau v0.6)
\einheitenkopf{Einheit 5 · Körper im Koordinatensystem und Drehungen}

\erg{1}{a) Sie liegt parallel zur $x_1x_2$-Ebene in der Höhe $4$.}
\erg{2}{a) $F$ liegt senkrecht über $C$ in der Höhe von $D$: $F(2 | 4 | 3)$ (nicht mit $E$ verwechseln). \quad b) $\overrightarrow{AD} = \overrightarrow{EH} = (0 | 3 | 0)$, $\overrightarrow{AE} = (0 | 0 | 5)$: $G = B + \overrightarrow{AD} + \overrightarrow{AE} = (6 | 4 | 5)$ (nicht $C$). \quad c) $\overrightarrow{AD} = \overrightarrow{EH} = (0 | 5 | 0)$, $\overrightarrow{AE} = (0 | 0 | 3)$: $G = B + \overrightarrow{AD} + \overrightarrow{AE} = (4 | 7 | 4)$ (nicht $C$). \quad d) $F$ liegt senkrecht über $C$ in der Höhe von $D$: $F(8 | 1 | 7)$ (nicht mit $E$ verwechseln). \quad e) Z. B. rechter Winkel im Ursprung, die Katheten auf die Achsen: $(0 | 0 | 0)$, $(6 | 0 | 0)$, $(0 | 8 | 0)$, Spitze $(0 | 0 | 5)$ – die Hypotenuse $10$ gehört nicht auf eine Achse. \quad f) $C(15 | 15 | 0)$, also $G(15 | 15 | 7)$ und $H(0 | 15 | 7)$. \quad g) Nur die Ecke $(0 | 2 | 2)$ passt: danach $(-1 | 1 | 3)$. \quad h) Nur die Ecke $(4 | 0 | 0)$ passt: danach $(2 | 1 | -1)$. \quad i) Nur die Ecke $(3 | 0 | 2)$ passt: danach $(1 | -3 | 1)$. \quad j) Alle Grundecken haben dieselbe dritte Koordinate: Die Grundfläche ist parallel zur $x_1x_2$-Ebene. Höhe $h = 7 - 2 = 5$ (nicht $7$). \quad k) Alle Grundecken haben dieselbe dritte Koordinate: Die Grundfläche ist parallel zur $x_1x_2$-Ebene. Höhe $h = 3 - (-1) = 4$ (nicht $3$). \quad l) Alle Grundecken haben dieselbe dritte Koordinate: Die Grundfläche ist parallel zur $x_1x_2$-Ebene. Höhe $h = 9 - 3 = 6$ (nicht $9$).}
\erg{3}{a) Strahlensatz am Querschnitt, von der Traufe aus: $d = \frac{2 - 0{,}5}{5{,}5 - 0{,}5} \cdot 4 = 1{,}2$ m (Höhendifferenz, nicht $\frac{2}{5{,}5}$). Nicht zählend: zwei Streifen, Anteil $\frac{2 \cdot 1{,}2}{8} = 30\,\%$. \quad b) Strahlensatz am Querschnitt, von der Traufe aus: $d = \frac{1{,}8 - 1}{5 - 1} \cdot 5 = 1$ m (Höhendifferenz, nicht $\frac{1{,}8}{5}$). Nicht zählend: zwei Streifen, Anteil $\frac{2 \cdot 1}{10} = 20\,\%$. \quad c) Strahlensatz am Querschnitt, von der Traufe aus: $d = \frac{1 - 0{,}4}{2{,}9 - 0{,}4} \cdot 5 = 1{,}2$ m (Höhendifferenz, nicht $\frac{1}{2{,}9}$). Nicht zählend: zwei Streifen, Anteil $\frac{2 \cdot 1{,}2}{10} = 24\,\%$. \quad d) $C$ läuft auf einem Kreis um den Lotfußpunkt $L(4 | 4 | 0)$ auf $AB$ (nicht um den Ursprung), Radius $|\overrightarrow{LC}| = \sqrt{18}$. In der $x_1x_2$-Ebene senkrecht zu $AB$: $L \pm 3 \cdot (1 | 1 | 0)$; mit $x_1 < 4$: $C'(1 | 1 | 0)$. \quad e) $C$ läuft auf einem Kreis um den Lotfußpunkt $L(2 | 5 | 0)$ auf $AB$ (nicht um den Ursprung), Radius $|\overrightarrow{LC}| = 5$. In der $x_1x_2$-Ebene senkrecht zu $AB$: $L \pm 5 \cdot (1 | 0 | 0)$; mit $x_1 > 2$: $C'(7 | 5 | 0)$. \quad f) Ecke $(k | k | k)$ auf $CS$: $C + t \cdot \overrightarrow{CS} = (4 - 4t | 4 - 4t | 6t)$, alle drei Koordinaten gleich: $4 - 4t = 6t$, also $t = 0{,}4$ und $k = 2{,}4$. Die Seitenflächen $SBC$ und $SCD$ liegen in $3x_1 + 2x_3 = 12$ und $3x_2 + 2x_3 = 12$; alle acht Würfelecken erfüllen dort „$\le$“ und liegen damit in der Pyramide – weil beide Körper konvex sind, liegt der ganze Würfel darin (die Prüfung einer Ecke allein genügt nicht). \quad g) Ecke $(k | k | k)$ auf $CS$: $C + t \cdot \overrightarrow{CS} = (3 - 3t | 3 - 3t | 6t)$, alle drei Koordinaten gleich: $3 - 3t = 6t$, also $t = \frac{1}{3}$ und $k = 2$. Die Seitenflächen $SBC$ und $SCD$ liegen in $2x_1 + x_3 = 6$ und $2x_2 + x_3 = 6$; alle acht Würfelecken erfüllen dort „$\le$“ und liegen damit in der Pyramide – weil beide Körper konvex sind, liegt der ganze Würfel darin (die Prüfung einer Ecke allein genügt nicht). \quad h) Für $2 \le s \le 5$ ein Dreieck (drei Ecken): die Bodenstrecke von $x_1 = 2$ bis $x_1 = 6$ und die Spitze auf dem First in der Höhe $3$ – immer dasselbe Dreieck, Flächeninhalt $\frac12 \cdot 4 \cdot 3 = 6$ (Grenzen sind die $x_2$-Koordinaten von $P$ und $Q$). Für $1 < s < 2$ und $5 < s < 7$ ein Trapez (vier Ecken): Die Dreiecksflächen $ABP$ bzw. $CDQ$ schneiden die Spitze ab. \quad i) Für $3 \le s \le 6$ ein Dreieck (drei Ecken): die Bodenstrecke von $x_1 = 0$ bis $x_1 = 6$ und die Spitze auf dem First in der Höhe $4$ – immer dasselbe Dreieck, Flächeninhalt $\frac12 \cdot 6 \cdot 4 = 12$ (Grenzen sind die $x_2$-Koordinaten von $P$ und $Q$). Für $1 < s < 3$ und $6 < s < 9$ ein Trapez (vier Ecken): Die Dreiecksflächen $ABP$ bzw. $CDQ$ schneiden die Spitze ab. \quad j) Die Ebene enthält die Deckendiagonale $EG$ und schneidet den Boden in der dazu parallelen Strecke von $S_t$ bis $(4 | 4 - t | 0)$. Für $0 \le t < 4$ ein Viereck (gleichschenkliges Trapez), für $t = 4$ ein gleichseitiges Dreieck (drei Achsen): $S_{4}$ ist die Ecke $B$, die beiden Bodenpunkte fallen zusammen (nicht dem Viereck zuschlagen). Genau zwei Symmetrieachsen nur für $t = 0$: dann ist es ein Rechteck. \quad k) Die Ebene enthält die Deckendiagonale $EG$ und schneidet den Boden in der dazu parallelen Strecke von $S_t$ bis $(6 | 6 - t | 0)$. Für $0 \le t < 6$ ein Viereck (gleichschenkliges Trapez), für $t = 6$ ein gleichschenkliges Dreieck: $S_{6}$ ist die Ecke $B$, die beiden Bodenpunkte fallen zusammen (nicht dem Viereck zuschlagen). Genau zwei Symmetrieachsen nur für $t = 0$: dann ist es ein Rechteck.}
